\documentclass[10pt, letterpaper]{article}

% Packages:
\usepackage[
  ignoreheadfoot, % set margins without considering header and footer
  top=2 cm, % seperation between body and page edge from the top
  bottom=2 cm, % seperation between body and page edge from the bottom
  left=1.7 cm, % seperation between body and page edge from the left
  right=1.7 cm, % seperation between body and page edge from the right
  footskip=1.0 cm, % seperation between body and footer
  % showframe % for debugging 
]{geometry} % for adjusting page geometry
\usepackage{titlesec} % for customizing section titles
\usepackage{tabularx} % for making tables with fixed width columns
\usepackage{array} % tabularx requires this
\usepackage[dvipsnames]{xcolor} % for coloring text
\definecolor{primaryColor}{RGB}{0, 0, 0} % define primary color
\usepackage{enumitem} % for customizing lists
\usepackage{fontawesome5} % for using icons
\usepackage{amsmath} % for math
\usepackage[
  pdftitle={Junseo Lee's CV},
  pdfauthor={Junseo Lee},
  pdfcreator={LaTeX with RenderCV},
  colorlinks=true,
  urlcolor=primaryColor
]{hyperref} % for links, metadata and bookmarks
\usepackage[pscoord]{eso-pic} % for floating text on the page
\usepackage{calc} % for calculating lengths
\usepackage{bookmark} % for bookmarks
\usepackage{lastpage} % for getting the total number of pages
\usepackage{changepage} % for one column entries (adjustwidth environment)
\usepackage{paracol} % for two and three column entries
\usepackage{ifthen} % for conditional statements
\usepackage{needspace} % for avoiding page brake right after the section title
\usepackage{iftex} % check if engine is pdflatex, xetex or luatex

% Ensure that generate pdf is machine readable/ATS parsable:
\ifPDFTeX
  \input{glyphtounicode}
  \pdfgentounicode=1
  \usepackage[T1]{fontenc}
  \usepackage[utf8]{inputenc}
  \usepackage{lmodern}
\fi

\usepackage{charter}
\linespread{1.12} % slightly looser line spacing
\clubpenalty=10000 \widowpenalty=10000 % never leave a single line of a paragraph alone across pages

% Some settings:
\raggedright
\AtBeginEnvironment{adjustwidth}{\partopsep0pt} % remove space before adjustwidth environment
\pagestyle{plain} % page number centered in the footer
\setcounter{secnumdepth}{0} % no section numbering
\setlength{\parindent}{0pt} % no indentation
\setlength{\topskip}{0pt} % no top skip
\setlength{\columnsep}{0.15cm} % set column seperation
%\pagenumbering{gobble} % no page numbering

\titleformat{\section}{\needspace{5\baselineskip}\bfseries\fontsize{13.5}{16}\selectfont}{}{0pt}{\MakeUppercase}[\vspace{1pt}\titlerule]

\titlespacing{\section}{
  % left space:
  -1pt
}{
  % top space:
  0.5 cm
}{
  % bottom space:
  0.25 cm
} % section title spacing

\renewcommand\labelitemi{$\vcenter{\hbox{\small$\bullet$}}$} % custom bullet points
\newenvironment{highlights}{
  \begin{itemize}[
    topsep=0.10 cm,
    parsep=0.10 cm,
    partopsep=0pt,
    itemsep=0.12 cm,
    leftmargin=0 cm + 10pt
  ]
}{
  \end{itemize}
} % new environment for highlights


\newenvironment{highlightsforbulletentries}{
  \begin{itemize}[
    topsep=0.10 cm,
    parsep=0.10 cm,
    partopsep=0pt,
    itemsep=0pt,
    leftmargin=10pt
  ]
}{
  \end{itemize}
} % new environment for highlights for bullet entries

\newenvironment{onecolentry}{
  \begin{adjustwidth}{
    0 cm + 0.00001 cm
  }{
    0 cm + 0.00001 cm
  }
}{
  \end{adjustwidth}
} % new environment for one column entries

\newenvironment{twocolentry}[2][]{
  \onecolentry
  \def\secondColumn{#2}
  \setcolumnwidth{\fill, 6 cm}
  \begin{paracol}{2}
}{
  \switchcolumn \raggedleft \secondColumn
  \end{paracol}
  \endonecolentry
} % new environment for two column entries

\newenvironment{threecolentry}[3][]{
  \onecolentry
  \def\thirdColumn{#3}
  \setcolumnwidth{, \fill, 4.5 cm}
  \begin{paracol}{3}
  {\raggedright #2} \switchcolumn
}{
  \switchcolumn \raggedleft \thirdColumn
  \end{paracol}
  \endonecolentry
} % new environment for three column entries

\newenvironment{header}{
  \setlength{\topsep}{0pt}\par\kern\topsep\centering\linespread{1.5}
}{
  \par\kern\topsep
} % new environment for the header

\newcommand{\placelastupdatedtext}{% \placetextbox{<horizontal pos>}{<vertical pos>}{<stuff>}
  \AddToShipoutPictureFG*{% Add <stuff> to current page foreground
  \put(
    \LenToUnit{\paperwidth-2 cm-0 cm+0.05cm},
    \LenToUnit{\paperheight-1.0 cm}
  ){\vtop{{\null}\makebox[0pt][c]{
    \small\color{gray}\textit{Last updated in November 2026}\hspace{\widthof{Last updated in November 2026}}
  }}}%
  }%
}%

% save the original href command in a new command:
\let\hrefWithoutArrow\href

% new command for external links:


% publication entry: venue label on the left, title/authors/venue on the right
\newcommand{\me}{\underline{\textbf{Junseo Lee}}}
\newcommand{\pubentry}[4]{%
  \begin{onecolentry}
  \begin{tabularx}{\linewidth}{@{}p{2.2cm}>{\raggedright\arraybackslash}X@{}}
    \textbf{#1} & #2 \\[1pt]
    & \mbox{\me}, #3 \\[1pt]
    & \textit{#4}
  \end{tabularx}
  \end{onecolentry}
}

\begin{document}
  \placelastupdatedtext
  \newcommand{\AND}{\unskip
    \cleaders\copy\ANDbox\hskip\wd\ANDbox
    \ignorespaces
  }
  \newsavebox\ANDbox
  \sbox\ANDbox{$|$}

  \begin{header}
    \fontsize{25 pt}{25 pt}\selectfont Junseo Lee

    \vspace{5 pt}

    \normalsize
    %\mbox{Seoul, Republic of Korea}%
    %\kern 5.0 pt%
    %\AND%
    %\kern 5.0 pt%
    \mbox{\hrefWithoutArrow{mailto:junseo.lee@snu.ac.kr}{junseo.lee@snu.ac.kr}}%
    \kern 5.0 pt%
    \AND%
    \kern 5.0 pt%
    %\mbox{\hrefWithoutArrow{tel:+82-10-9414-4887}{+82-10-9414-4887}}%
    %\kern 5.0 pt%
    %\AND%
    %\kern 5.0 pt%
    \mbox{\hrefWithoutArrow{https://junseolee0.github.io/}{junseolee0.github.io}}%
    \kern 5.0 pt%
    \AND%
    \kern 5.0 pt%
    %\mbox{\hrefWithoutArrow{https://www.linkedin.com/in/junseo-lee-144795301/}{linkedin.com/in/junseo-lee-144795301}}%
    \mbox{\hrefWithoutArrow{https://scholar.google.com/citations?user=f2Bx8o8AAAAJ&hl=en}{Google Scholar}}%
    %\kern 5.0 pt%
    %\AND%
    %\kern 5.0 pt%
    %\mbox{\hrefWithoutArrow{https://github.com/yourusername}{github.com/yourusername}}%
  \end{header}

  \vspace{5 pt - 0.3 cm}


  \fontsize{10.5}{13}\selectfont % body text size (both pages)
  \section{Research Interests}
%    \begin{onecolentry}
%      My research interests are in \textbf{computer architecture} and
%      \textbf{the architectural support for emerging workloads},
%      particularly vision/graphics tasks such as neural rendering.
%    \end{onecolentry}
%
%    \vspace{0.2 cm}
%
    \begin{onecolentry}
      Computer Architecture, GPU Microarchitecture, Hardware-Software Co-Design for Emerging Workloads, Computer Vision/Graphics, Neural Rendering
    \end{onecolentry}

  \section{Education}
    
    \begin{twocolentry}{
      Seoul, Republic of Korea
    }
      \textbf{Seoul National University}
    \end{twocolentry}
    \vspace{0.1cm}
    \begin{twocolentry}{
      Mar 2022 – Feb 2028 (expected)
    }
      M.S./Ph.D., Electrical and Computer Engineering
    \end{twocolentry}
    
    \begin{onecolentry}
      \begin{highlightsforbulletentries}
      \item \hrefWithoutArrow{https://comparch.snu.ac.kr}{Computer Architecture and Systems Lab} (Advisor: \hrefWithoutArrow{https://jaewoong.org}{Prof. Jaewoong Sim})
      \end{highlightsforbulletentries}
    \end{onecolentry}
    
    \vspace{0.1cm}
    \begin{twocolentry}{
      Mar 2018 – Feb 2022
    }
      B.S., Electrical and Computer Engineering
    \end{twocolentry}
    
  \section{Publications}
    \needspace{4\baselineskip}
    \pubentry{\makebox[0pt][l]{\raisebox{-1.2\baselineskip}[0pt][0pt]{\normalfont\small\textit{to appear}}}HPCA '27}{BlendCore: Architectural Support for High-Throughput Programmable Blending}{\mbox{Wonjae Jang}, \mbox{Sangyun Jeon}, \mbox{Jaewoong Sim}}{33rd International Symposium on High-Performance Computer Architecture, 2027}
    \vspace{0.40 cm}
    \needspace{4\baselineskip}
    \pubentry{HPCA '26}{GRTX: Efficient Ray Tracing for 3D Gaussian-Based Rendering}{\mbox{Sangyun Jeon}, \mbox{Jungi Lee}, \mbox{Junyong Park}, \mbox{Jaewoong Sim}}{32nd International Symposium on High-Performance Computer Architecture, 2026}
    \vspace{0.40 cm}
    \needspace{4\baselineskip}
    \pubentry{HPCA '25}{VR-Pipe: Streamlining Hardware Graphics Pipeline for Volume Rendering}{\mbox{Jaisung Kim}, \mbox{Junyong Park}, \mbox{Jaewoong Sim}}{31st International Symposium on High-Performance Computer Architecture, 2025}
    \vspace{0.40 cm}
    \needspace{4\baselineskip}
    \pubentry{ASPLOS '24}{GSCore: Efficient Radiance Field Rendering via Architectural Support for 3D Gaussian Splatting}{\mbox{Seokwon Lee}, \mbox{Jungi Lee}, \mbox{Junyong Park}, \mbox{Jaewoong Sim}}{2024 International Conference on Architectural Support for Programming Languages and Operating Systems, 2024}
    \vspace{0.40 cm}
    \needspace{4\baselineskip}
    \pubentry{ISCA '23}{NeuRex: A Case for Neural Rendering Acceleration}{\mbox{Kwanseok Choi}, \mbox{Jungi Lee}, \mbox{Seokwon Lee}, \mbox{Joonho Whangbo}, \mbox{Jaewoong Sim}}{50th Annual International Symposium on Computer Architecture, 2023}

  \section{Honors and Awards}
    \begin{onecolentry}
      \textbf{Outstanding Graduate Student Research Fellowship}: Department of ECE, Seoul National University \hfill 2026
    \end{onecolentry}
    \vspace{0.15 cm}
    \begin{onecolentry}
      \textbf{Student Travel Grants}: ISCA '23, ASPLOS '24, HPCA '25 \hfill 2023--2025
    \end{onecolentry}

  \clearpage
  \fontsize{10.5}{13}\selectfont
  \section{Research Experience}
    \begin{twocolentry}{
      Mar 2022 – Present
    }
      \textbf{Graduate Research Assistant}
    \end{twocolentry}
%
    \begin{twocolentry}{
      Seoul, Republic of Korea
    }
      Seoul National University (Advisor: \hrefWithoutArrow{https://jaewoong.org}{Prof. Jaewoong Sim})
    \end{twocolentry}

    \vspace{0.10 cm}
    \begin{onecolentry}
      \begin{highlights}
%        \item Advisor: \hrefWithoutArrow{https://jaewoong.org}{Prof. Jaewoong Sim}
        \item \textbf{BlendCore (HPCA '27): Extending the graphics pipeline in GPUs for high-throughput shader blending}
        \begin{itemize}
            \item As a follow-up to VR-Pipe, offloaded blending from the fixed-function ROP to high-throughput SMs to resolve the blending bottleneck.
            \item Designed a lightweight tile distribution unit that dynamically steers blending work to lightly loaded SMs with only a few kilobytes of storage per Graphics Processing Cluster (GPC).
        \end{itemize}
        \item \textbf{GRTX (HPCA '26): Software and hardware optimizations for efficient ray tracing of 3D Gaussians}
        \begin{itemize}
            \item Proposed a streamlined acceleration structure in which all Gaussians share a single unit-sphere geometry with per-Gaussian transformations to substantially reduce BVH size and traversal overhead.
            \item Designed hardware support for traversal checkpointing in ray tracing units to eliminate redundant node visits during multi-round tracing.
        \end{itemize}
        \item \textbf{VR-Pipe (HPCA '25): A hardware extension to the graphics pipeline in GPUs for volume rendering}
        \begin{itemize}
            \item Extended the hardware graphics pipeline to reduce the pressure on the fixed-function blending unit (ROP).
            \item Designed hardware support for early termination to skip unnecessary fragments and quad merging to pre-blend fragments in shader cores before the ROP.
        \end{itemize}
        \item \textbf{GSCore (ASPLOS '24): A specialized accelerator for the 3D Gaussian Splatting rendering pipeline}
        \begin{itemize}
            \item Proposed three algorithmic optimizations and tailored hardware to reduce ineffective computations.
            \item Showed that the algorithmic optimizations alone speed up rendering on GPUs in software.
        \end{itemize}
        \item \textbf{NeuRex (ISCA '23): Algorithm-hardware co-design for Neural Radiance Fields (NeRF) rendering}
        \begin{itemize}
            \item Made the multi-resolution hash encoding hardware-friendly and co-designed an encoding engine, which together eliminate irregular off-chip memory accesses.
            \item Implemented NeuRex in SystemVerilog with small area overhead and low power consumption.
        \end{itemize}
      \end{highlights}
    \end{onecolentry}
  
%  \section{Publications}
%    \begin{samepage}
%      \needspace{4\baselineskip}
%      \begin{onecolentry}
%        \textbf{BlendCore: Architectural Support for High-Throughput Programmable Blending} \\
%        \mbox{\textbf{Junseo Lee}}, \mbox{Wonjae Jang}, \mbox{Sangyun Jeon}, \mbox{Jaewoong Sim} \\
%        \textit{Proc. of the 33rd International Symposium on High-Performance Computer Architecture (\textbf{HPCA}), March 2027 (to appear)}
%      \end{onecolentry}
%      \vspace{0.30 cm}
%      \needspace{4\baselineskip}
%      \begin{onecolentry}
%        \textbf{GRTX: Efficient Ray Tracing for 3D Gaussian-Based Rendering} \\
%        \mbox{\textbf{Junseo Lee}}, \mbox{Sangyun Jeon}, \mbox{Jungi Lee}, \mbox{Junyong Park}, \mbox{Jaewoong Sim} \\
%        \textit{Proc. of the 32nd International Symposium on High-Performance Computer Architecture (\textbf{HPCA}), January 2026}
%      \end{onecolentry}
%      \vspace{0.30 cm}
%      \needspace{4\baselineskip}
%      \begin{onecolentry}
%        \textbf{VR-Pipe: Streamlining Hardware Graphics Pipeline for Volume Rendering} \\
%        \mbox{\textbf{Junseo Lee}}, \mbox{Jaisung Kim}, \mbox{Junyong Park}, \mbox{Jaewoong Sim} \\
%        \textit{Proc. of the 31st International Symposium on High-Performance Computer Architecture (\textbf{HPCA}), March 2025}
%      \end{onecolentry}
%      \vspace{0.30 cm}
%      \needspace{4\baselineskip}
%      \begin{onecolentry}
%        \textbf{GSCore: Efficient Radiance Field Rendering via Architectural Support for 3D Gaussian Splatting} \\
%        \mbox{\textbf{Junseo Lee}}, \mbox{Seokwon Lee}, \mbox{Jungi Lee}, \mbox{Junyong Park}, \mbox{Jaewoong Sim} \\
%        \textit{Proc. of the 2024 International Conference on Architectural Support for Programming Languages and Operating Systems (\textbf{ASPLOS}), April 2024}
%      \end{onecolentry}
%      \vspace{0.30 cm}
%      \needspace{4\baselineskip}
%      \begin{onecolentry}
%        \textbf{NeuRex: A Case for Neural Rendering Acceleration} \\
%        \mbox{\textbf{Junseo Lee}}, \mbox{Kwanseok Choi}, \mbox{Jungi Lee}, \mbox{Seokwon Lee}, \mbox{Joonho Whangbo}, \mbox{Jaewoong Sim} \\
%        \textit{Proc. of the 50th Annual International Symposium on Computer Architecture (\textbf{ISCA}), June 2023}
%      \end{onecolentry}
%    \end{samepage}
    
%  \section{Projects}
%    \begin{twocolentry}{
%      \href{https://github.com/sinaatalay/rendercv}{github.com/name/repo}
%    }
%      \textbf{Multi-User Drawing Tool}\end{twocolentry}
%
%    \vspace{0.10 cm}
%    \begin{onecolentry}
%      \begin{highlights}
%        \item Developed an electronic classroom where multiple users can simultaneously view and draw on a "chalkboard" with each person's edits synchronized
%        \item Tools Used: C++, MFC
%      \end{highlights}
%    \end{onecolentry}
%
%    \vspace{0.2 cm}
%
%    \begin{twocolentry}{
%      \href{https://github.com/sinaatalay/rendercv}{github.com/name/repo}
%    }
%      \textbf{Synchronized Desktop Calendar}\end{twocolentry}
%
%    \vspace{0.10 cm}
%    \begin{onecolentry}
%      \begin{highlights}
%        \item Developed a desktop calendar with globally shared and synchronized calendars, allowing users to schedule meetings with other users
%        \item Tools Used: C\#, .NET, SQL, XML
%      \end{highlights}
%    \end{onecolentry}
%
%    \vspace{0.2 cm}
%
%    \begin{twocolentry}{
%      2002
%    }
%      \textbf{Custom Operating System}\end{twocolentry}
%
%    \vspace{0.10 cm}
%    \begin{onecolentry}
%      \begin{highlights}
%        \item Built a UNIX-style OS with a scheduler, file system, text editor, and calculator
%        \item Tools Used: C
%      \end{highlights}
%    \end{onecolentry}

  \section{Teaching}
    \begin{onecolentry}
      \textbf{Graduate Teaching Assistant}
    \end{onecolentry}
%
    \begin{twocolentry}{
      Seoul, Republic of Korea
    }
      Seoul National University
    \end{twocolentry}
    \vspace{0.1cm}
    \begin{onecolentry}
      \begin{highlightsforbulletentries}
        \item ECE 315.A: Digital Systems Design and Experiments \hfill Fall 2022
        \item ECE 322: Computer Organization \hfill Spring 2022
      \end{highlightsforbulletentries}
    \end{onecolentry}
    \vspace{0.1cm}
    \begin{onecolentry}
      Led lab sessions and the final project that implements a CNN accelerator in RTL on an FPGA (ECE 315.A). \\
      Prepared assignments, including a C++-based cycle-level cache simulator (ECE 322).
    \end{onecolentry}
  
  \section{Skills}
    
    \begin{onecolentry}
      \textbf{Languages:} C/C++, CUDA, Python, Verilog/SystemVerilog, GLSL
    \end{onecolentry}
    \vspace{0.2 cm}
    \begin{onecolentry}
      \textbf{Tools:} PyTorch, architectural simulators (e.g., gem5, GPGPU-Sim, MacSim, Ramulator), Intel VTune
    \end{onecolentry}
    \vspace{0.2 cm}
    \begin{onecolentry}
      \textbf{APIs:} OpenGL, Vulkan, OptiX
    \end{onecolentry}
\end{document}